% ===========================================================================
\chapter{Architecture de la pipeline DevSecOps}
\label{chap:archi-pipeline}
% ===========================================================================

\begin{constat}{état de ce chapitre}
Ce chapitre présente la conception d'ensemble de la pipeline et la
correspondance entre pratiques, outils et étages. Les étages de la phase~0 sont
en service ; les étages d'analyse de sécurité sont mis en œuvre en phase~5 et
documentés au chapitre~\ref{chap:stages}.
\end{constat}

\section{Principe de conception}

La pipeline est structurée par la cartographie des pratiques sur les phases du
cycle de vie DevSecOps établie par Prates et Pereira~\cite{prates2025}. Cette
cartographie répond à la question de l'emplacement : elle indique à quel moment
de la chaîne chaque pratique doit intervenir.

\begin{table}[H]
\centering
\caption{Correspondance entre phases du cycle DevSecOps, pratiques et étages de
notre pipeline}
\label{tab:cycle-pipeline}
\small
\begin{tabular}{@{}llp{7.0cm}@{}}
\toprule
\textbf{Phase du cycle} & \textbf{Pratiques} & \textbf{Mise en œuvre} \\
\midrule
Create
& \pid{P8} \pid{P9} \pid{P20}
& Crochets \emph{pre-commit}, \texttt{CODEOWNERS}, guide de codage sécurisé \\
\addlinespace
Plan
& \pid{P7} \pid{P23} \pid{P24}
& Modélisation des menaces, analyse de risques, périmètre ASVS \\
\addlinespace
Adapt
& \pid{P6} \pid{P12}
& Intégration continue, infrastructure en tant que code \\
\addlinespace
Verify
& \pid{P1} \pid{P2} \pid{P3} \pid{P4} \pid{P11}
& SAST, DAST, suite de tests d'abus, analyse de vulnérabilités, revue \\
\addlinespace
Preproduction
& \pid{P4} \pid{P16} \pid{P22}
& Analyse d'image, test d'intrusion, durcissement \\
\addlinespace
Release
& \pid{P10}
& Déploiement continu conditionné à la réussite des contrôles \\
\addlinespace
Predict / Detect
& \pid{P5} \pid{P18}
& Supervision continue, journalisation, alertes \\
\addlinespace
Respond
& \pid{P28}
& Procédure de gestion des incidents \\
\addlinespace
Configure
& \pid{P30} \pid{P21}
& Gestion des secrets, gestion de configuration \\
\bottomrule
\end{tabular}
\end{table}

\captureaplacer{figures/architecture-pipeline.png}{Schéma de la pipeline
DevSecOps : étages, déclencheurs et points de blocage, superposés aux phases du
cycle de vie. À produire en phase 5, lorsque tous les étages seront en service.}

\section{Étages prévus et outils retenus}

\subsection{Règle de traçabilité des outils}

Un principe a été appliqué à l'ensemble des choix d'outillage.

\begin{decision}{outils issus de la source scientifique}
À l'exception de deux outils explicitement signalés, tous les outils retenus
figurent dans le tableau~8 de Prates et Pereira~\cite{prates2025}. Chaque choix
est donc rattachable à la source de référence du module.

Les écarts sont signalés comme tels et justifiés en propre. Cette distinction
est maintenue dans tout le rapport : elle évite de présenter une préférence
personnelle comme une recommandation de la littérature.
\end{decision}

\subsection{Tableau des étages}

\begin{table}[H]
\centering
\caption{Étages de sécurité de la pipeline}
\label{tab:etages}
\small
\begin{tabular}{@{}llp{3.5cm}p{4.3cm}@{}}
\toprule
\textbf{Étage} & \textbf{Prat.} & \textbf{Outil} & \textbf{Déclencheur} \\
\midrule
Détection de secrets & \pid{P30} & Gitleaks & chaque envoi, historique complet \\
Hygiène du dépôt     & \pid{P8}  & pre-commit & chaque envoi \\
SAST                 & \pid{P1}  & SonarQube, ESLint-security & envoi et \emph{pull request} \\
Composition logicielle & \pid{P13} & OWASP Dependency-Check & \emph{pull request} et nocturne \\
Dépendances JS       & \pid{P4}  & Retire.js & \emph{pull request} \\
Image de conteneur   & \pid{P4}  & Trivy & après construction d'image \\
Infrastructure       & \pid{P4}  & Checkov & modification de manifeste \\
Tests d'abus         & \pid{P3}  & Jest, Supertest & \emph{pull request} \\
DAST                 & \pid{P2}  & OWASP ZAP & référence en PR, complet en nocturne \\
Conformité continue  & \pid{P17} & Checkov, politique & \emph{pull request} \\
\bottomrule
\end{tabular}
\end{table}

Retire.js mérite un commentaire particulier : cet outil est spécifique à
l'écosystème JavaScript, ce qui constitue une adéquation directe au périmètre du
Projet~3.

\section{Politique de blocage}

\begin{constat}{la politique de blocage est elle-même une décision de sécurité}
Déterminer quels résultats font échouer la construction et lesquels ne
produisent qu'un avertissement n'est pas un réglage technique mineur.

Faire échouer la construction sur chaque résultat de faible gravité paralyse le
travail et enseigne à l'équipe le contournement de la pipeline. Ne rien faire
échouer rend la pipeline décorative. La position intermédiaire doit être
justifiée explicitement, et non supposée.
\end{constat}

La politique retenue sera consignée dans le registre ADR~0008 en phase~5, et
reproduite au chapitre~\ref{chap:stages}. Elle précisera pour chaque outil le
seuil de gravité bloquant, ainsi que la procédure d'acceptation d'un résultat
--- par qui, sur quel motif, et avec quelle échéance de réexamen.

\section{Séparation entre contrôles rapides et analyses profondes}

Deux contraintes s'opposent. Les contrôles doivent être suffisamment rapides
pour s'exécuter à chaque \emph{pull request} sans entraver le travail ; les
analyses dynamiques complètes ne peuvent pas l'être.

La conception retenue sépare donc les déclencheurs : analyse ZAP de référence et
analyse de composition rapide sur chaque \emph{pull request}, analyse ZAP
complète et analyse de composition approfondie sur déclenchement nocturne
planifié.

Ce choix est à la fois une décision de sécurité et une décision de performance.
Son effet mesuré est présenté au chapitre~\ref{chap:performance}.
